\documentclass{article}

%\usepackage[binary-units]{siunitx}[=v2]
\usepackage{siunitx}

\usepackage{xcolor}
\usepackage{cancel}
\usepackage{multirow}
\usepackage{longtable}
%\usepackage{sansmath} % not in LaTeXML

\usepackage{booktabs}
\usepackage{etoolbox}

\usepackage{hyperref}

\DeclareSIQualifier\polymer{pol}
\DeclareSIQualifier\catalyst{cat}
\DeclareSIQualifier\isotropic{i} %?
\newbox\slashbox
\setbox\slashbox\hbox{\verb|\|}
\def\Backslash{\copy\slashbox}

\DeclareSIUnit\abbra{ a }
\DeclareSIUnit\abbrb{ \abbra }  % abbreviation
\DeclareSIUnit\abbrc{ \abbrb }  % abbreviation twice
\DeclareSIUnit\abbrd{ \abbre }  % before defined
\DeclareSIUnit\abbre{ e }

\def\mabbra{a}
\DeclareSIUnit\mabbrb{\mabbra}

\DeclareSIPrefix\killer{k}{3}
% so six_get_prefix_from_power(3) -> \killer, not \kilo
% but they have the same abbreviation, so it works out

\DeclareSIUnit\noop{\relax}

\DeclareSIPower\quartic\tothefourth{4}

\begin{document}

{\sloppy
\tableofcontents
}

\section{Introduction}

This tries to follow the package documentation released 2025-07-09, with additional tests.  Items remaining (generally ordered from hardest to easiest) are
\begin{itemize}
\item error: evaluate-expression, expression  % 4 occurrences
\item sansmath: package needed to evaluate two more tests
\item The spacing in tables is not quite right.  See the comment on \&six\_table\_space
\item font and spacing: several instances where we should make a font change or adjust the spacing, but don't.  % 10 of each
\item detect if loaded with \verb|[=v2]| or \verb|[=2021-04-09]| or earlier  % see issue 2719
\end{itemize}

\section{siunitx for the impatient}
\num{12345,67890}\\
\num{.3e45}\\
\complexnum{1+-2i}\\
\numproduct{1.654 x 2.34 x 3.430}

\unit{kg.m.s^{-1}} \\
\unit{\kilogram\metre\per\second} \\
\unit[per-mode = symbol]{\kilogram\metre\per\second} \\
\unit[per-mode = symbol]{\kilogram\metre\per\ampere\per\second}

\numlist{10;20;30} \\
\qtylist{0.13;0.67;0.80}{\milli\metre} \\
\numrange{10}{20} \\
\qtyrange{0.13}{0.67}{\milli\metre}

\subsection*{Other tests}
%\ang{1.2e10;2 x 3;4}
\numlist{10 ; 20 ; 30} \\
\ang{1.2;3;4}

Unsemantic:
\unit{m^2.s}\\
\unit{\mu m^2}

Semantic again:
\qty{0.094}{\pi . \milli\meter . \milli\radian}\\
\qty{0.094}{\frac{1}{3} . \milli\meter . \milli\radian}\\
\qty{0.094}{\pi \per \milli\meter . \milli\radian\tothe{3}}

\section{Using the siunitx package}

\subsection{Numbers}

\num{123}\\
\num{1234}\\
\num{12345}\\
\num{0.123}\\
\num{0,1234}\\
\num{.12345}\\
\num{3.45d-4}\\
\num{-e10}

\numlist{10;30;50;70}

\numproduct{10 x 30}

\numrange{10}{30}

\subsubsection{Other Tests}

\num{123e4}\\
%\num{123e4(3)}\\
\num{123(3)e4}

\num{123\pm2}\\
\complexnum{123\pm2i}\\
\complexnum{123+234i}\\
\complexnum{123(1)+234(1)ie3}\\
\complexnum{234(1)ie3}
%\num{123+234}\\
%\num{123e2+234e3i}\\

\def\dig{1234}
\def\odd{\xi}
\def\odder{\odd}

Pretty nonsensical stuff?
\num[parse-numbers=false]{1.\pi e+3}\\
\num{\dig.\dig}\\
\num[input-digits=\xi,parse-numbers=false]{3\xi}\\
\num[input-digits=\xi,parse-numbers=false]{3\odd}\\
%\num[input-digits=\odd]{3\odd}\\
\num[input-digits=\odd, parse-numbers=false]{3\odd}\\
\num[input-digits=\odd, parse-numbers=false]{3\odder}\\
\num[input-digits=\odder, parse-numbers=false]{3\odder}

\num{1.23(1)}\\
%\num{1.23(0.01)}
\num{1.23\pm0.01}\\
%\num{1.23(1.0)}\\
%\num{1.23(\pm1)}\\
\num[parse-numbers=false]{1.23(\pi)}

\num[parse-numbers=false]{1.234(5) x \pi} \\
\num[uncertainty-mode = separate, parse-numbers=false]{1.234(5) x \pi}

\subsubsection{Lists and ranges of numbers}

\qtylist{10;30;45}{\metre}\\
\qtyproduct{10 x 30 x 45}{\metre}\\
\qtyrange{10}{30}{\metre}

\subsection{Angles}

\ang{10} \\
\ang{12.3} \\
\ang{4,5} \\
\ang{1;2;3} \\
\ang{;;1} \\
\ang{+10;;} \\
\ang{-0;1;}

\subsection{Units}

\subsubsection{unit}
\unit{kg.m/s^2} \\
\unit{g_{polymer}~mol_{cat}.s^{-1}} \\
\unit{\kilo\gram\metre\per\square\second} \\
\unit{\gram\per\cubic\centi\metre} \\
\unit{\square\volt\cubic\lumen\per\farad} \\
\unit{\metre\squared\per\gray\cubic\lux} \\
\unit{\henry\second}

\subsubsection{qty}
\qty[mode = text]{1.23}{J.mol^{-1}.K^{-1}} \\
\qty{.23e7}{\candela} \\
\qty[per-mode = symbol]{1.99}{\per\kilogram} \\
\qty[per-mode = fraction]{1,345}{\coulomb\per\mole}

\subsubsection{qtylist}
\qtylist{10;30;45}{\metre}

\subsubsection{qtyproduct}
\qtyproduct{10 x 30 x 45}{\metre}

\subsubsection{qtyrange}
\qtyrange{10}{30}{\metre}


\subsection{Complex numbers and quantities}

\complexnum{123+234ie3}\\
\complexnum{123(1)+234(1)ie3}\\
%\num{123e2(1)+234e3i(1)}\\
%\num{+234(1)ie3}\\              %comes out weird?
\complexnum{+3i}\\
\complexnum{+3ie4}

\subsection{The unit macros}
\newcommand{\showunit}[1]{#1 & \texttt{\Backslash#1} & \unit{\csname#1\endcsname}}

The SI base units:
\begin{center}
\begin{tabular}{lll}\toprule
Unit & Macro & Symbol \\\midrule
\showunit{ampere}\\
\showunit{candela}\\
\showunit{kelvin}\\
\showunit{kilogram}\\
\showunit{metre}\\
\showunit{mole}\\
\showunit{second}\\\bottomrule
\end{tabular}
\end{center}

The coherent derived units:
\begin{center}
\begin{tabular}{@{}llc llc@{}}\toprule
Unit & Macro & Symb &
Unit & Macro & Symb \\\cmidrule(r){1-3}\cmidrule(l){4-6}
\showunit{becquerel}     & \showunit{newton} \\
\showunit{degreeCelsius} & \showunit{ohm} \\
\showunit{coulomb}       & \showunit{pascal} \\
\showunit{farad}         & \showunit{radian} \\
\showunit{gray}          & \showunit{siemens} \\
\showunit{hertz}         & \showunit{sievert} \\
\showunit{henry}         & \showunit{steradian} \\
\showunit{joule}         & \showunit{tesla} \\
\showunit{katal}         & \showunit{volt} \\
\showunit{lumen}         & \showunit{watt} \\
\showunit{lux}           & \showunit{weber} \\\bottomrule
\end{tabular}
\end{center}

The non-SI units:
\begin{center}
\begin{tabular}{lll}\toprule
Unit & Macro & Symbol \\\midrule
\showunit{astronomicalunit}\\
\showunit{bel}\\
\showunit{dalton}\\
\showunit{day}\\
\showunit{decibel}\\
\showunit{degree}\\
\showunit{electronvolt}\\
\showunit{hectare}\\
\showunit{hour} \\
\showunit{litre}\\
\showunit{liter}\\
\showunit{arcminute}\\
\showunit{minute}\\
\showunit{arcsecond}\\
\showunit{neper}\\
\showunit{tonne}\\\bottomrule
\end{tabular}
\end{center}

%\ExplSyntaxOn
% this accesses siunitx internals, and may therefore break
\newcommand{\showprefix}[1]{
  #1 & \texttt{\Backslash#1} & \unit{ \csname #1\endcsname \noop } &
  \qty[prefix-mode=extract-exponent, print-unity-mantissa=false]{1}{\csname#1\endcsname\noop}%
% the following
%  \str_if_eq:nnTF { #1 } { micro } { -6 } { \__ltxml_power_from_text:n { #1 } }
% works in LaTeX, but not LaTeXML (b/c it doesn't have the internals)
}
%\cs_new:Npn \__ltxml_prefix_from_text:n #1 {
%  \use:c { __siunitx_unit_ \token_to_str:c { #1 } :w }
%}
%\cs_new:Npn \__ltxml_power_from_prefix:n #1 {
%  \prop_get:NeNTF
%  \l__siunitx_unit_prefixes_forward_prop
%  { #1 }
%  \l_tmpa_tl
%  { \l_tmpa_tl }
%  { \qty [ prefix-mode=extract-exponent, print-unity-mantissa=false ] { 1 } { \use:c { #1 } \noop } }
%}
%\cs_new:Npn \__ltxml_power_from_text:n #1 {
%  {
%    \bool_set_false:N \l__siunitx_unit_test_bool
%    \bool_set_false:N \l__siunitx_unit_parsing_bool
%    \exp_args:Ne \__ltxml_power_from_prefix:n { \__ltxml_prefix_from_text:n { #1 } }
%  }
%}
%\__ltxml_prefix_from_text:n { kibi }
%\ExplSyntaxOff

SI prefixes:
\begin{center}
\begin{tabular}{@{}llcl llcl@{}}\toprule
Prefix & Cmd & Symb & Power &
Prefix & Cmd & Symb & Power \\\cmidrule(r){1-4}\cmidrule(l){5-8}
\showprefix{quecto} & \showprefix{deca}\\
\showprefix{ronto}  & \showprefix{hecto}\\
\showprefix{yocto}  & \showprefix{kilo}\\
\showprefix{zepto}  & \showprefix{mega}\\
\showprefix{atto}   & \showprefix{giga}\\
\showprefix{femto}  & \showprefix{tera}\\
\showprefix{pico}   & \showprefix{peta}\\
\showprefix{nano}   & \showprefix{exa}\\
\showprefix{micro}  & \showprefix{zetta}\\
\showprefix{milli}  & \showprefix{yotta}\\
\showprefix{centi}  & \showprefix{ronna}\\
\showprefix{deci}   & \showprefix{quetta}\\\bottomrule
\end{tabular}
\end{center}

\unit{\square\becquerel} \\
\unit{\joule\squared\per\lumen} \\
\unit{\cubic\lux\volt\tesla\cubed} \\
\unit{\henry\tothe{5}} \\
\unit{\raiseto{4.5}\radian} \\
\unit{\joule\per\mole\per\kelvin} \\
\unit{\joule\per\mole\kelvin} \\
\unit{\per\henry\tothe{5}} \\
\unit{\per\square\becquerel} \\
\unit{\kilogram\of{metal}} \\
\qty[qualifier-mode = bracket]{0.1}{\milli\mole\of{cat}\per\kilogram\of{prod}} \\
\unit[per-mode = fraction]{\cancel\kilogram\metre\per\cancel\kilogram\per\second} \\
\unit{\highlight{red}\kilogram\metre\per\second} \\
\unit[unit-color = purple]{\highlight{blue}\kilogram\metre\per\second} % todo font

\subsection{Unit abbreviations}

\newcommand{\showunitabbrev}[4]{#1#3 & \texttt{\Backslash#2#4} & \unit{\csname#2#4\endcsname}}

\ExplSyntaxOn
\newcommand{\showunitabbrevs}[3]{
 \midrule
 \clist_map_variable:nNn { #3 } { \l_tmpa_tl } {
  \exp_after:wN \showunitabbrev \l_tmpa_tl { #1 } { #2 } \\
 }
}
\ExplSyntaxOff

Abbreviated units

\begin{center}
\begin{tabular}{lcc}\toprule
Unit & Abbreviation & Symbol \\\midrule
\showunitabbrevs{gram}{g}{{femto}{f},{pico}{p},{nano}{n},{micro}{u},{milli}{m},{}{},{kilo}{k}}
\bottomrule
\end{tabular}

\begin{tabular}{lcc}\toprule
Unit & Abbreviation & Symbol \\\midrule
\showunitabbrevs{metre}{m}{{pico}{p},{nano}{n},{micro}{u},{milli}{m},{centi}{c},{deci}{d},{}{},{kilo}{k}}
\bottomrule
\end{tabular}

\begin{tabular}{lcc}\toprule
Unit & Abbreviation & Symbol \\\midrule
\showunitabbrevs{meter}{m}{{pico}{p},{nano}{n},{micro}{u},{milli}{m},{centi}{c},{deci}{d},{}{},{kilo}{k}}
\bottomrule
\end{tabular}

\begin{tabular}{lcc}\toprule
Unit & Abbreviation & Symbol \\\midrule
\showunitabbrevs{second}{s}{{atto}{a},{femto}{f},{pico}{p},{nano}{n},{micro}{u},{milli}{m},{}{}}
\bottomrule
\end{tabular}

\begin{tabular}{lcc}\toprule
Unit & Abbreviation & Symbol \\\midrule
\showunitabbrevs{mole}{mol}{{femto}{f},{pico}{p},{nano}{n},{micro}{u},{milli}{m},{}{},{kilo}{k}}
\bottomrule
\end{tabular}

\begin{tabular}{lcc}\toprule
Unit & Abbreviation & Symbol \\\midrule
\showunitabbrevs{ampere}{A}{{pico}{p},{nano}{n},{micro}{u},{milli}{m},{}{},{kilo}{k}}
\bottomrule
\end{tabular}

\begin{tabular}{lcc}\toprule
Unit & Abbreviation & Symbol \\\midrule
\showunitabbrevs{litre}{l}{{micro}{u},{milli}{m},{}{},{hecto}{h}}
\bottomrule
\end{tabular}

\begin{tabular}{lcc}\toprule
Unit & Abbreviation & Symbol \\\midrule
\showunitabbrevs{liter}{L}{{micro}{u},{milli}{m},{}{},{hecto}{h}}
\bottomrule
\end{tabular}

\begin{tabular}{lcc}\toprule
Unit & Abbreviation & Symbol \\\midrule
\showunitabbrevs{hertz}{Hz}{{milli}{m},{}{},{kilo}{k},{mega}{M},{giga}{G},{tera}{T}}
\bottomrule
\end{tabular}

\begin{tabular}{lcc}\toprule
Unit & Abbreviation & Symbol \\\midrule
\showunitabbrevs{newton}{N}{{milli}{m},{}{},{kilo}{k},{mega}{M}}
\bottomrule
\end{tabular}

\begin{tabular}{lcc}\toprule
Unit & Abbreviation & Symbol \\\midrule
\showunitabbrevs{pascal}{Pa}{{}{},{kilo}{k},{mega}{M},{giga}{G}}
\bottomrule
\end{tabular}

\begin{tabular}{lcc}\toprule
Unit & Abbreviation & Symbol \\\midrule
\showunitabbrevs{ohm}{ohm}{{milli}{m},{kilo}{k},{mega}{M}}
\bottomrule
\end{tabular}

\begin{tabular}{lcc}\toprule
Unit & Abbreviation & Symbol \\\midrule
\showunitabbrevs{volt}{V}{{pico}{p},{nano}{n},{micro}{u},{milli}{m},{}{},{kilo}{k}}
\bottomrule
\end{tabular}

\begin{tabular}{lcc}\toprule
Unit & Abbreviation & Symbol \\\midrule
\showunitabbrevs{watt}{W}{{nano}{n},{micro}{u},{milli}{m},{}{},{kilo}{k},{mega}{M},{giga}{G}}
\bottomrule
\end{tabular}

\begin{tabular}{lcc}\toprule
Unit & Abbreviation & Symbol \\\midrule
\showunitabbrevs{joule}{J}{{micro}{u},{milli}{m},{}{},{kilo}{k}}
\bottomrule
\end{tabular}

\begin{tabular}{lcc}\toprule
Unit & Abbreviation & Symbol \\\midrule
\showunitabbrevs{electronvolt}{eV}{{milli}{m},{}{},{kilo}{k},{mega}{M},{giga}{G},{tera}{T}}
\bottomrule
\end{tabular}

\begin{tabular}{lcc}\toprule
Unit & Abbreviation & Symbol \\\midrule
\showunitabbrevs{watt hour}{Wh}{{kilo}{k}}
\bottomrule
\end{tabular}

\begin{tabular}{lcc}\toprule
Unit & Abbreviation & Symbol \\\midrule
\showunitabbrevs{farad}{F}{{femto}{f},{pico}{p},{nano}{n},{micro}{u},{milli}{m},{}{}}
\bottomrule
\end{tabular}

\begin{tabular}{lcc}\toprule
Unit & Abbreviation & Symbol \\\midrule
\showunitabbrevs{henry}{H}{{femto}{f},{pico}{p},{nano}{n},{micro}{u},{milli}{m},{}{}}
\bottomrule
\end{tabular}

\begin{tabular}{lcc}\toprule
Unit & Abbreviation & Symbol \\\midrule
\showunitabbrevs{coulomb}{C}{{nano}{n},{micro}{u},{milli}{m},{}{}}
\bottomrule
\end{tabular}

\begin{tabular}{lcc}\toprule
Unit & Abbreviation & Symbol \\\midrule
\showunitabbrevs{tesla}{T}{{micro}{u},{milli}{m},{}{}}
\bottomrule
\end{tabular}

\begin{tabular}{lcc}\toprule
Unit & Abbreviation & Symbol \\\midrule
\showunitabbrevs{kelvin}{K}{{}{}}
\bottomrule
\end{tabular}

\begin{tabular}{lcc}\toprule
Unit & Abbreviation & Symbol \\\midrule
\showunitabbrevs{bel}{B}{{deci}{d}}
\bottomrule
\end{tabular}
\end{center}

Binary prefixes:
\newcounter{tablerow}
\begin{center}
% the documentation numbers are hardcoded in
% /usr/local/texlive/YYYY/texmf-dist/source/latex/siunitx/siunitx.tex
% the powers are not stored in siunitx.  see siunitx-binary.dtx
\newcommand{\showBinaryPrefix}[1]{%
  \stepcounter{tablerow}%
  #1 & \texttt{\Backslash#1} & \unit{\csname #1\endcsname\noop} & \arabic{tablerow}0
}
\begin{tabular}{cccc}\toprule
Prefix & Command & Symbol & Power \\\midrule
\showBinaryPrefix{kibi} \\
\showBinaryPrefix{mebi} \\
\showBinaryPrefix{gibi} \\
\showBinaryPrefix{tebi} \\
\showBinaryPrefix{pebi} \\
\showBinaryPrefix{exbi} \\
\showBinaryPrefix{zebi} \\
\showBinaryPrefix{yobi} \\\bottomrule
\end{tabular}
\end{center}

\subsection{Creating new macros}

\qty{67890}{\degree} \\
\qty[quantity-product = \,]{67890}{\degree}

\unit{\kilogram\tothefourth}\\
\unit{\quartic\metre}

\qty{1.234}{\gram\polymer\per\mole\catalyst\per\hour}

\subsubsection{Other Tests}
% Note that \kilogram is defined as \kilo\gram, and acts sorta like a macro
\unit{\kilogram}
\unit{\kilo\gram}
% Not allowed.
%\unit{\kilo\kilogram}
%\unit{\kilo\kilo\gram}
{
% and note that regular macros ARE expanded!
\def\killermeater{\kilo\meter}
\unit{\killermeater}
}
{
% BUT, it prefers its OWN definitions!
\def\gram{Grr}
\unit{\kilo\gram}
}

% However, it's happy with various indirect definitions.
% Ie. units, prefixes (at least) act like regular macros within \si.
\unit{\abbra}\\
\unit{\abbrb}\\
\unit{\abbrc}\\
\unit{\abbrd}\\
\unit{\abbre}

\unit{\mabbra}\\
\unit{\mabbrb}

\unit{\killer\meter}

% Where can highlight go?
\unit{\highlight{red}\kilo\gram\metre\per\second} \\
\unit{\kilo\highlight{red}\gram\metre\per\second} \\
\unit{\kilo\gram\highlight{red}\metre\per\second} \\
\unit{\kilo\gram\metre\highlight{red}\per\second} \\
\unit{\kilo\gram\metre\per\highlight{red}\second}

\unit{\cancel\kilo\gram\metre\per\second} \\
\unit{\kilo\cancel\gram\metre\per\second} \\
\unit{\kilo\gram\cancel\metre\per\second} \\
\unit{\kilo\gram\metre\cancel\per\second} \\
\unit{\kilo\gram\metre\per\cancel\second}


\subsection{Tabular material}

Standard behaviour of the \texttt{S} column type
\begin{center}
\begin{tabular}{S}
\toprule
{Some Values} \\
\midrule
2.3456 \\
34.2345 \\
-6.7835 \\
90.473 \\
5642.5 \\
1.2e3 \\
e4 \\
\bottomrule
\end{tabular}
\end{center}

Detection of surrounding material in an \texttt{S} column
\begin{center}
\begin{tabular}{S[color=orange]}
\toprule
{Some Values} \\
\midrule
12.34 \\
\color{purple} 975,31 \\
44.268 \textsuperscript{\emph{a}} \\
\bottomrule
\end{tabular}
\end{center}

Text before and after numbers
\begin{center}
\sisetup{table-format = {now }2.4{\textsuperscript{\emph{a}}}}
\begin{tabular}{@{}S@{}}
\toprule
{Values} \\
\midrule
2.3456 \\
34.2345 \textsuperscript{\emph{a}}\\
34.2345 {\textsuperscript{\emph{a}}}\\
56.7835 \\
%now~ 90.473 \\  % todo error
{now~} 90.473 \\  % todo incorrect
\bottomrule
\end{tabular}
\end{center}

Controlling complex alignment with the tablenum macro
\begin{center}
\begin{tabular}{lr}
\toprule
Heading & Heading \\
\midrule
Info & More info \\
Info & More info \\
\multicolumn{2}{c}{\tablenum[table-format = 4.4]{12,34}} \\
\multicolumn{2}{c}{\tablenum[table-format = 4.4]{333.5567}} \\
\multicolumn{2}{c}{\tablenum[table-format = 4.4]{4563.21}} \\
\bottomrule
\end{tabular}
\qquad
\begin{tabular}{lr}
\toprule
Heading & Heading \\
\midrule
\multirow{2}*{\tablenum{88,999}} & aaa \\
& bbb \\
\multirow{2}*{\tablenum{33,435}} & ccc \\
& ddd \\
\bottomrule
\end{tabular}
\end{center}

\section{Package control options}

\stepcounter{subsection}

\subsection{Printing}
\subsubsection{mode, number-mode, unit-mode}
\subsubsection{reset-text-[family\textbar series\textbar shape]} % todo font
{
 \sisetup{mode = text}
 {\itshape \num{1234}}\\
 {\bfseries \num{1234}}\\
 {\sffamily \num{1234}}\\
 \sisetup{
 reset-text-family = false ,
 reset-text-series = false ,
 reset-text-shape = false
 }
 {\itshape \num{1234}}\\
 {\bfseries \num{1234}}\\
 {\sffamily \num{1234}}
}

\subsubsection{propagate-math-font, reset-math-version}
{
 {\boldmath \unit{\kilogram}}\\ % todo font
 %{\sansmath $\unit{\kilogram}$}\\ % todo usepackage sansmath
 {$\mathsf{\unit{\kilogram}}$}\\
 \sisetup{
 propagate-math-font = true ,
 reset-math-version = false
 }
 {\boldmath \unit{\kilogram}}\\
 %{\sansmath $\unit{\kilogram}$}\\ % todo usepackage sansmath
 {$\mathsf{\unit{\kilogram}}$}
}

\subsubsection{text-[family\textbar series]-to-math} % todo font
{
 {\sffamily \unit{\kilogram}}\\
 {\bfseries $\unit{\kilogram}$}\\
 \sisetup{
 text-family-to-math = true ,
 text-series-to-math = true
 }
 {\sffamily \unit{\kilogram}}\\
 {\bfseries $\unit{\kilogram}$}
}

\subsubsection{text-font-command}
{
 \sisetup{number-mode = text}
 \qty{123456789}{\kilo\volt\per\centi\metre} \\
 \sisetup{text-font-command = \fontfamily{pplx}\selectfont}
 \qty{123456789}{\kilo\volt\per\centi\metre}  % todo font
}

\subsubsection{text-[sub\textbar super]script-command}
{
 \sisetup{unit-mode = text}
 \unit{\kg\of{polymer}} \\
 \newcommand*\mysubscript[1]{\textsubscript{\textcolor{blue}{#1}}}
 \newcommand\mysuperscript[1]{\textsuperscript{\textcolor{green}{#1}}}
 \sisetup{text-subscript-command = \mysubscript}
 \unit{\kg\of{polymer}} % todo font
}

\subsubsection{color, number-color, unit-color} % todo font
{
 \color{red}%
 Some text \\
 \qty{4}{\kilogram} \\
 More text \\
 \qty[color = blue]{4}{\kilogram} \\
 \qty[number-color = blue]{4}{\kilogram} \\
 \qty[unit-color = blue]{4}{\kilogram} \\
 Still red here!
}

\qty[number-color=blue,color=green]{1}{\gram} \\
\qty[color=green,number-color=blue]{1}{\gram} \\
\qty[unit-color=blue,color=green]{1}{\gram} \\
\qty[color=green,unit-color=blue]{1}{\gram} \\
\ang[number-color=green,unit-color=blue]{-1;2;3} \\
\numlist{1;2;3}\\
\numproduct[product-mode=phrase]{1x2}\\
\numrange{1}{2}

\numproduct{1.234(5) x 6.78(9)}

\subsection{Parsing numbers}

%\subsubsection{input-digits, input-decimal-markers, input-signs, input-exponent-markers}
\subsubsection{input-digits, input-ignore}
\num[input-ignore=3a\pi]{1a234\pi5}

\subsubsection{input-comparators}

\num{< 10} \\
\qty{>> 5}{\metre} \\
\num{\le 0.12}

\subsubsection{input-[open\textbar close]-uncertainty, input-uncertainty-signs}
\num{9.99(9)}\\
\num{9.99 +- 0.09}\\
\num{9.99 \pm 0.09}\\
\num{123 +- 4.5}\\
\num{12.3 +- 6}\\
\num{123.4(12)} \\
\num{123.4(1.2)}\\
\num{123.4(12)(45)} \\
\num{123.4 \pm 1.2 \pm 4.5}

\subsubsection{input-uncertainty-divider}
\num{10.56(12:34)} \\
\num{123.4(4.6:7.8)}\\
\num{10.56(1:2)(3)}\\
\num{6.45(2)(3:4)}

Also:\\
\num{5.6(1.2:3.4)} \\
\num{78.56(12:34)} \\
\num{123.4(12)(45)(67)} \\
\num{123.4 \pm 1.2 \pm 4.5 \pm 6.7}


\subsubsection{parse-numbers}
\num[parse-numbers = false]{\sqrt{2}}\\
\qty[parse-numbers = false]{\sqrt{3}}{\metre}

\subsubsection*{Other tests}
\num[parse-numbers=false]{\pi}\\
\num[parse-numbers=false]{2\pi}\\
\num[parse-numbers=false]{\pi/3}

% todo error
%\subsubsection{evaluate-expression, expression}
%{
%\sisetup{evaluate-expression}%
%\qty{2 + 4 * 3}{\joule} \\
%\qty[expression = 10 * (#1)]{2 + 4 * 3}{\joule}
%}

\subsubsection{retain-explicit-decimal-marker, retain-explicit-plus, retain-negative-zero, retain-zero-uncertainty}
%\num[retain-explicit-decimal-marker]{10.} \\  % v2 or v3
\num{+345} \\
\num[retain-explicit-plus]{+345} \\
\num{-0} \\
\num[retain-negative-zero]{-0} \\
\num{12.3(0)} \\
\num[retain-zero-uncertainty]{12.3(0)}

{
\num{456} \\
\num{.789} \\
\sisetup{
  add-decimal-zero = false,
  add-integer-zero = false,
}%
\num{456} \\
\num{.789} % v2: .789; v3: 0.789. But add-integer-zero is v2
}

\subsection{Post-processing numbers}

\subsubsection{exponent-mode, fixed-exponent}
{
\num{0.001} \\
\num{0.0100} \\
\num{1200} \\
\sisetup{exponent-mode = scientific}%
\num{0.001} \\
\num{0.0100} \\
\num{1200} \\
\sisetup{exponent-mode = engineering}%
\num{0.001} \\
\num{0.0100} \\
\num{1200} \\
\sisetup{
exponent-mode = fixed,
fixed-exponent = 2,
}%
\num{0.001} \\
\num{0.0100} \\
\num{1200} \\
}
\num{1.23e4} \\
\num[exponent-mode = fixed, fixed-exponent = 0]{1.23e4}

\complexnum[exponent-mode=engineering]{0.5+1000i}
% todo tex has (0.000 500+1.000i)x10^3 (see siunitx#813)

\subsubsection{exponent-thresholds}
This is unique to version 3.

\subsubsection{drop-exponent, drop-uncertainty}
\num{0.01(2)} \\
\num[drop-uncertainty]{0.01(2)} \\
\num{0.01e3} \\
\num[drop-exponent]{0.01e3}  % causes warning: Potentially ambiguous dropping of exponent

\subsubsection{round-mode, round-precision, round-pad}

{
\num{1.23456} \\
\num{14.23} \\
\num{0.12345(9)} \\
\num{0.0012345}\\
\sisetup{
  round-mode = places,
  round-precision = 3
}%
places:\\
\num{1.23456} \\
\num{14.23} \\
\num{0.12345(9)} \\
\num{0.0012345}\\
\sisetup{
  round-mode = figures,
  round-precision = 3
}%
figures:\\
\num{1.23456} \\
\num{14.23} \\
\num{0.12345(9)} \\
\num{9999}\\
\num{0.0012345}\\
\sisetup{
round-mode = uncertainty,
round-precision = 1
}%
uncertainty:\\
\num{0.12345(9)} \\
\num{0.12345(23)} \\
\num{0.12345(234)}\\
\num{12345(123)}\\
\num{123.456(012)}\\
\num{1234.5(12.3)}\\
\num{1234.56(98)}\\
\num{9.99(99)}\\
\num{99.99(999)}\\
}%
{
\sisetup{round-mode = figures, round-precision = 4}%
figures:\\
\num{12.3} \\
\num{012.34} \\
\num[round-pad = false]{12.3}\\
}
Also\\
\numproduct[round-mode=places]{.123 x .456 x .789}

\subsubsection{round-direction, round-half}
{
\sisetup{round-mode = places}
\num{0.054} \\
\num{0.046} \\
\sisetup{round-direction = down}%
\num{0.054} \\
\num{0.046} \\
\sisetup{round-direction = up}%
\num{0.054} \\
\num{0.046} \\
}
{
\sisetup{
round-mode = figures,
round-precision = 1,
round-half = up
}%
\num{0.055} \\
\num{0.045} \\
\sisetup{round-half = even}%
\num{0.055} \\
\num{0.045}
}

\subsubsection{uncertainty-round-direction}
{
\sisetup{round-mode = uncertainty}
\num{0.123(41)} \\
\sisetup{uncertainty-round-direction = up}%
\num{0.123(41)} \\ % lets try again
\sisetup{round-precision=1}%
\num{0.123(41)}
}

\subsubsection{round-minimum}
{
\sisetup{round-mode = places}%
\num{0.0055} \\
\num{0.0045} \\
\num{-0.0045} \\
\sisetup{round-minimum = 0.01}%
\num{0.0055} \\
\num{0.0045} \\
\num{-0.0045}\\
\sisetup{round-minimum = {0,01}}%
\num{0,0045}
}

\subsubsection{round-zero-positive}
{
\sisetup{round-mode = places}%
\num{-0.001} \\
\sisetup{round-zero-positive = false}%
\num{-0.001}
}

\subsubsection{drop-zero-decimal}
\num{2.0} \\
\num{2.1} \\
{
\sisetup{drop-zero-decimal}%
\num{2.0} \\
\num{2.1}
}

\subsubsection{minimum-decimal-digits, minimum-integer-digits}
\num{123} \\
\num[minimum-integer-digits = 2]{123} \\
\num[minimum-integer-digits = 4]{123} \\
\num{0.123} \\
\num[minimum-decimal-digits = 2]{0.123} \\
\num[minimum-decimal-digits = 4]{0.123}

\subsubsection{Other Tests}
{
\num{001} \\
%\num{123.} \\  % v2 or v3
\num{456} \\
\num{.789} \\
\sisetup{
  minimum-decimal-digits = 0,
  print-zero-integer = false,
}%
%\num{123.} \\  % v2 or v3
\num{456} \\
\num{.789} \\
}
Also: \\
\num[minimum-decimal-digits=2]{1} \\
\num[minimum-decimal-digits=1, drop-zero-decimal]{1.00} \\
\num[minimum-decimal-digits=2, drop-zero-decimal]{1}

% an example from https://github.com/josephwright/siunitx/issues/756
\num[separate-uncertainty=true, round-mode=uncertainty, drop-zero-decimal=true, round-precision=1]{-52.01(43)}


\subsection{Printing Numbers}
\subsubsection{group-digits, group-separator}
\num{12345.67890} \\
\num[group-digits = none]{12345.67890} \\
\num[group-digits = decimal]{12345.67890} \\
\num[group-digits = integer]{12345.67890}\\
\num{12345} \\
\num[group-separator = {,}]{12345} \\
\num[group-separator = \text{~}]{12345}\\
\num[group-separator = \ ]{12345}

\subsubsection{group-minimum-digits}
\num{1234} \\
\num{12345} \\
\num[group-minimum-digits = 5]{1234} \\
\num[group-minimum-digits = 5]{12345} \\
\num{1234.5678} \\
\num{12345.67890} \\
\num[group-minimum-digits = 5]{1234.5678} \\
\num[group-minimum-digits = 5]{12345.67890}

\subsubsection*{Other tests}
\num{1234} \\
\num[group-minimum-digits = 4]{1234} \\
\num{1234.5678} \\
\num[group-minimum-digits = 4]{1234.5678}

\subsubsection{digit-group-size, digit-group-first-size, digit-group-other-size}
\num{1234567890} \\
\num[digit-group-size = 5]{1234567890} \\
\num[digit-group-other-size = 2]{1234567890}

\subsubsection*{Other tests}
\num{1234567890.1234567890}\\
\num[digit-group-size=4]{1234567890.1234567890}\\
\num[digit-group-size = 5, digit-group-other-size = 2]{1234567890} \\
\num[digit-group-other-size = 2, digit-group-size = 5]{1234567890} \\
\num[digit-group-size = 5, digit-group-first-size = 2]{1234567890} \\
\num[digit-group-first-size = 2, digit-group-size = 5]{1234567890}

\subsubsection{output-decimal-marker}
\num{1.23} \\
\num[output-decimal-marker = {,}]{1.23} % todo spacing: the , gets interpreted as punctuation

\subsubsection{exponent-base, exponent-product}
\num[exponent-product = \times]{1e2} \\
\num[exponent-product = \cdot]{1e2} \\
\num[exponent-base = 2]{1e2}

\subsubsection{output-exponent-marker}
\num[output-exponent-marker = e]{1e2} \\
\num[output-exponent-marker = \text{e}]{1e2} \\
\num[output-exponent-marker = \mathrm{E}]{1e2}\\
\num[output-exponent-marker = \ensuremath{\mathrm{E}}]{1e2}

\subsubsection{uncertainty-[mode\textbar separator], output-[open\textbar close]-uncertainty}
{
\num{123.45(120)} \\
\num{0.035(14)} \\
\sisetup{uncertainty-mode = full}
\num{123.45(120)} \\
\num{0.035(14)} \\
\sisetup{uncertainty-mode = compact-marker}
\num{123.45(120)} \\
\num{0.035(14)} \\
%\sisetup{uncertainty-mode = separate}
\sisetup{
output-open-uncertainty = [,
output-close-uncertainty = ],
uncertainty-separator = \,
}%
\num{1.234(5)} \\
\num{1.234\pm 0.005} \\
\num[uncertainty-mode = separate]{1.234(5)} \\
\num[uncertainty-mode = separate]{1.234\pm 0.005} \\
\num{8.2(13)} \\
\num{8.2\pm1.3} \\
\num[uncertainty-mode = separate]{8.2(13)}\\
\num[uncertainty-mode = separate]{8.2\pm1.3} \\
\num{1.2 +- 0.001}\\
\num[uncertainty-mode = separate]{1.2 +- 0.001}\\
}
\fbox{\parbox{4em}{%
\sisetup{uncertainty-mode = separate}%
\num{67890+-12345}\\
\num[allow-uncertainty-breaks=false]{67890+-12345}% todo spacing
}}

\subsubsection{uncertainty-descriptors, uncertainty-descriptor-mode, uncertainty-descriptor-separator} % todo spacing, todo font
{
\num{1.2(3)(4)} \\
\sisetup{uncertainty-descriptors = {sys, stat}}
\num{1.2(3)(4)} \\ % default uncertainty-descriptor-mode = bracket-separator
\num[uncertainty-descriptor-mode = subscript]{1.2(3)(4)}\\
\num[uncertainty-descriptor-mode = bracket]{1.2(3)(4)}\\
\num[uncertainty-descriptor-mode = separator]{1.2(3)(4)}\\
\num[uncertainty-mode=separate]{1.2(3)}\\
\num[uncertainty-mode=separate]{1.2(3)(4)(5)}\\
}
\num[uncertainty-mode=separate, uncertainty-descriptors={sys}]{1.2(3)}\\
\num[uncertainty-mode=separate, uncertainty-descriptors={sys,stat,third}]{1.2(3)(4)}

\subsubsection{simplify-uncertainty}
\num{10.56(2:2)} \\
\num[simplify-uncertainty]{10.56(2:2)}

\subsubsection{bracket-ambiguous-numbers}
\num{1 e10} \\
\complexnum{2i e10} \\
\complexnum{1+2i e10} \\
\complexnum[bracket-ambiguous-numbers = false]{1+2i e10} \\
{
\sisetup{uncertainty-mode = separate}
\num{1.2(3)e4} \\
\num[bracket-ambiguous-numbers = false]{1.2(3)e4} \\
}
{
\sisetup{uncertainty-mode = separate, bracket-ambiguous-numbers = false}%
\num{1.234(5)e-4} %\\
%\qty{1.234(5)e-4}{\metre}
% "This option only applies to pure numbers: when formatting quantities,
% the need for brackets also depends on the placement of units, so is controlled by separate-uncertainty-units."
}

\subsubsection{negative-color}
\num{-15673} \\
\num[negative-color = red]{-15673}

also:\\
\num[negative-color = red]{-0}\\
\num[negative-color = red, retain-negative-zero]{-0}\\
\num[color=green,negative-color=red]{-10}\\
\num[negative-color=red,color=green]{-10}

{\sisetup{negative-color=red}
\complexnum{1}\\
\complexnum{-1}\\
\complexnum{3i}\\
\complexnum{-3i}\\
\complexnum{1+3i}\\
\complexnum{-1+3i}\\
\complexnum{1-3i}\\
\complexnum{-1-3i}
}

\subsubsection{bracket-negative-numbers}
\num{-15673} \\
\num[bracket-negative-numbers]{-15673} \\
\qty{-10}{\metre} \\
\qty[bracket-negative-numbers]{-10}{\metre}

{\sisetup{bracket-negative-numbers}
\num{<<-5}\\
\num{-13(1)}\\
\num{-13+-1}\\
\num{-e10}\\
\num{-2e10}\\
\complexnum{1}\\
\complexnum{-1}\\
\complexnum{3i}\\
\complexnum{-3i}\\
\complexnum{1+3i}\\
\complexnum{-1+3i}\\
\complexnum{1-3i}\\
\complexnum{-1-3i}
}

\subsubsection{tight-spacing} % todo spacing
\num{2e3} \\
\num[tight-spacing = true]{2e3}

\subsubsection{print-implicit-plus, print-[exponent\textbar mantissa]-implicit-plus}
a\\
\num{345} \\
\num[print-implicit-plus]{345} \\
\num[print-mantissa-implicit-plus]{345} \\
\num[print-exponent-implicit-plus]{345} \\
b\\
\num{1e2} \\
\num[print-implicit-plus]{1e2} \\
\num[print-mantissa-implicit-plus]{1e2} \\
\num[print-exponent-implicit-plus]{1e2} \\
c\\
\num{1e+2} \\
\num[print-implicit-plus]{1e+2} \\
\num[print-mantissa-implicit-plus]{1e+2} \\
\num[print-exponent-implicit-plus]{1e+2} \\
d\\
\num{-1e2} \\
\num[print-implicit-plus]{-1e2} \\
\num[print-mantissa-implicit-plus]{-1e2} \\
\num[print-exponent-implicit-plus]{-1e2} \\
e\\
\num{1e-2} \\
\num[print-implicit-plus]{1e-2} \\
\num[print-mantissa-implicit-plus]{1e-2} \\
\num[print-exponent-implicit-plus]{1e-2}

\subsubsection{print-unity-mantissa, print-zero-[exponent\textbar integer]}
\num{1e4} \\
\num[print-unity-mantissa = false]{1e4} \\
\num{444e0} \\
\num[print-zero-exponent = true]{444e0} \\
\num{0.123} \\
\num[print-zero-integer = false]{0.123}

\num{e4}\\
%\num[print-unity-mantissa]{e4}\\
\num[print-unity-mantissa=false]{e4}

\num[]{1e0}\\
\num[print-unity-mantissa=false]{1e0}\\ % 1 should leak through as per documentation b/c print-zero-exponent=false
\num[print-zero-exponent=true]{1e0}\\
\num[print-unity-mantissa=false, print-zero-exponent=true]{1e0}

\subsubsection{zero-decimal-as-symbol, zero-symbol}
\num{123.00} \\
{
\sisetup{zero-decimal-as-symbol}
\num{123.00} \\
%\num{123.} \\  % v2 or v3
\num{123} \\
\num[zero-symbol = \text{[{---}]}]{123.00}
}

\subsection{Lists, products and ranges}
\subsubsection{list-[final\textbar pair]-separator, list-separator}
\numlist{0.1;0.2;0.3}\\
\numlist[list-separator = {; }]{0.1;0.2;0.3}\\
\numlist[list-final-separator = {, }]{0.1;0.2;0.3} \\
\numlist[
list-separator = { and },
list-final-separator = { and finally }
]{0.1;0.2;0.3} \\
\numlist{0.1;0.2} \\
\numlist[list-pair-separator = {, and }]{0.1;0.2} \\
\numlist[color=blue]{0.1;0.2;0.3}

\subsubsection{product-[mode\textbar phrase\textbar symbol]}
\numproduct{5 x 100 x 2} \\
\numproduct[product-symbol = \ensuremath{\cdot}]{5 x 100 x 2} \\
{
\sisetup{product-mode = phrase}%
\numproduct{5 x 100 x 2}\\
\numproduct[product-phrase = { BY }]{5 x 100 x 2}
}

\subsubsection{range-open-phrase, range-phrase}
\numrange{5}{100} \\
\numrange[range-phrase = --]{5}{100}\\
\numrange{10}{12} \\
\numrange[range-open-phrase = {\text{from} }]{5}{100}\\
\numrange[color=blue,range-phrase = --]{5}{100}

\subsubsection{[list\textbar product\textbar range]-exponents}
\numlist{5e3;7e3;9e3;1e4} \\
\numproduct{5e3 x 7e3 x 9e3 x 1e4} \\
\numrange{5e3}{7e3} \\
{
\sisetup
{
list-exponents = combine-bracket ,
product-exponents = combine-bracket ,
range-exponents = combine-bracket
}
\numlist{5e3;7e3;9e3;1e4} \\
\numproduct{5e3 x 7e3 x 9e3 x 1e4} \\
\numrange{5e3}{7e3} \\
\numrange[range-open-phrase={\text{from }}]{5e3}{7e3} \\
\sisetup
{
list-exponents = combine ,
product-exponents = combine ,
range-exponents = combine
}
\numlist{5e3;7e3;9e3;1e4} \\
\numproduct{5e3 x 7e3 x 9e3 x 1e4} \\
\numrange{5e3}{7e3}
}

\subsubsection{[list\textbar product\textbar range]-units}
\qtylist{2;4;6;8}{\tesla} \\
\qtylist[list-units = bracket]{2;4;6;8}{\tesla} \\
\qtylist[list-units = repeat]{2;4;6;8}{\tesla} \\
\qtylist[list-units = single]{2;4;6;8}{\tesla} \\
\qtyrange{2}{4}{\degreeCelsius} \\
\qtyrange[range-units = bracket]{2}{4}{\degreeCelsius} \\
\qtyrange[range-units = repeat]{2}{4}{\degreeCelsius} \\
\qtyrange[range-units = single]{2}{4}{\degreeCelsius}\\
\qtyproduct{2 x 3 x 4}{\metre} \\
\qtyproduct[product-units = bracket]{2 x 3 x 4}{\metre}\\
\qtyproduct[product-units = repeat]{2 x 3 x 4}{\metre}\\
\qtyproduct[product-units = single]{2 x 3 x 4}{\metre}\\
\qtyproduct[product-units = bracket-power]{2 x 3 x 4}{\metre}\\
\qtyproduct[product-units = power]{2 x 3 x 4}{\metre}

\subsubsection{[list\textbar product\textbar range]-[close\textbar open]-bracket}
{
\sisetup{
list-units = bracket ,
list-open-bracket = [ ,
list-close-bracket = ]
}
\qtylist{5e3;7e3;9e3;1e4}{\second} \\
\sisetup{
product-units = bracket ,
product-open-bracket = \{ ,
product-close-bracket = \}
}
\qtyproduct{5e3 x 7e3 x 9e3 x 1e4}{\metre} \\
\sisetup{
range-units = bracket ,
range-open-bracket = \langle,
range-close-bracket = \rangle
}
\qtyrange{2}{4}{\degreeCelsius}\\
\sisetup{
list-units = bracket ,
list-open-bracket = ( ,
list-close-bracket = )
}
\numlist{5e3;7e3;9e3;1e4} \\
\sisetup{
product-units = bracket ,
product-open-bracket = ( ,
product-close-bracket = )
}
\numproduct{5e3 x 7e3 x 9e3 x 1e4} \\
\sisetup{
range-units = bracket ,
range-open-bracket = ( ,
range-close-bracket = )
}
\qtyrange{2}{4}{\degreeCelsius}
}

\subsubsection{[list\textbar product\textbar range]-independent-prefix}
{
a\\
\qtyrange{1e3}{1e9}{\watt}\\
\qtylist{1e3;1e9}{\watt}\\
\qtyproduct{1e3 x 1e9}{\watt}\\
b\\
\qtyrange[range-independent-prefix]{1e3}{1e9}{\watt}\\
\qtylist[list-independent-prefix]{1e3;1e9}{\watt}\\
$\qtyproduct[product-independent-prefix]{1e3 x 1e9}{\watt}$\\
\sisetup{exponent-to-prefix}%
c\\
\qtyrange{1000}{1e9}{\watt}\\
\qtyrange{1e3}{1e9}{\watt}\\
\qtylist{1000;1e9}{\watt}\\
\qtylist{1e3;1e9}{\watt}\\
\qtyproduct{1000 x 1e9}{\watt}\\
\qtyproduct{1e3 x 1e9}{\watt}\\
d\\
\qtyrange[range-independent-prefix]{1e3}{1e9}{\watt}\\
\qtylist[list-independent-prefix]{1e3;1e9}{\watt}\\
$\qtyproduct[product-independent-prefix]{1e3 x 1e9}{\watt}$
}


\subsubsection{prefix-mode = combine-exponent}
\qty{1700}{\g} \\
\qty{1.7e3}{\g} \\
{
\sisetup{prefix-mode = combine-exponent}%
\qty{1700}{\g} \\
\qty{1.7e3}{\g} \\
}
{
\sisetup{fixed-exponent = 3, exponent-mode = fixed}%
\qty{1700}{\g} \\
\qty{1.7e3}{\g}
}


\subsection{Complex Numbers}
\subsubsection{complex-mode}
\complexnum{1 + i} \\
\complexnum{1:45} \\
\complexnum[complex-mode = cartesian]{1 + i} \\
\complexnum[complex-mode = cartesian, round-mode = places]{1:45} \\
\complexnum[complex-mode = polar]{1 + i} \\
\complexnum[complex-mode = polar]{1:45}

\subsubsection{input-complex-roots}
\complexnum{9.99 + 88.8i} \\
\complexnum{9.99 + i88.8}

\subsubsection{output-complex-root} % todo font
\complexnum{1+2i} \\
\complexnum[output-complex-root = i]{1+2i} \\
\complexnum[output-complex-root = \text{\ensuremath{i}}]{1+2i} \\
\complexnum[output-complex-root = j]{1+2i}

\subsubsection{complex-root-position}
\complexnum{67-0.9i} \\
\complexnum[complex-root-position = before-number]{67-0.9i} \\
\complexnum[complex-root-position = after-number]{67-0.9i}

\subsubsection{complex-angle-unit, complex-symbol-angle, complex-symbol-degree}
\complexqty{1:1}{\ohm} \\
\complexqty[complex-angle-unit = radians]{1:1}{\ohm} \\
\complexqty[complex-symbol-degree = d]{1:1}{\ohm}
%\complexqty[complex-phase-command = \phase]{1:1}{\ohm} % needs steinmetz package

\subsubsection{print-complex-unity}
\complexqty{1i}{\ohm} \\
\complexqty[print-complex-unity]{1i}{\ohm}\\
\complexqty[print-complex-unity]{i}{\ohm}

\subsubsection{Other tests}
Real numbers as complex:\\
\complexnum{1}\\
\complexnum[print-complex-unity]{1}\\
\complexnum[complex-mode=cartesian]{1:180}\\
\complexnum[complex-mode=cartesian]{1:360}\\
\complexnum[complex-mode=cartesian]{1:540}\\
\complexnum[complex-mode=cartesian]{1:720}\\
\complexnum[complex-mode=cartesian, complex-angle-unit=radians]{1:3.14159265359}\\  % Perl uses sci notation for these two
\complexnum[complex-mode=cartesian, complex-angle-unit=radians]{1:6.28318530718}  % TeX doesn't

\complexnum[complex-mode=polar, complex-angle-unit=radians]{i}

\subsection{Angles}
\subsubsection{angle-mode}
\ang{2.67} \\
\ang{2;3;4} \\
\ang[angle-mode = arc]{2.67} \\
\ang[angle-mode = arc]{2;3;4} \\
\ang[angle-mode = decimal]{2.67} \\
\ang[angle-mode = decimal]{2;3;4}

\subsubsection{angle-separator}
\ang{6;7;6.5} \\ % todo spacing
\ang[angle-separator = \,]{6;7;6.5}

\subsubsection{fill-angle-[degrees\textbar minutes\textbar seconds]}
\ang{-1;;} \\
\ang{;-2;} \\
\ang{;;-3} \\
{
\sisetup{fill-angle-degrees}
\ang{-1;;} \\
\ang{;-2;} \\
\ang{;;-3} \\
}
{
\sisetup{fill-angle-minutes}
\ang{-1;;} \\
\ang{;-2;} \\
\ang{;;-3} \\
}
{
\sisetup{fill-angle-seconds}
\ang{-1;;} \\
\ang{;-2;} \\
\ang{;;-3} \\
}
\ang[fill-angle-minutes,fill-angle-seconds]{45.697}\\
\ang[color=blue,fill-angle-minutes,fill-angle-seconds]{45.697}

\subsubsection{angle-symbol-[degree\textbar minute\textbar second]}
{
\ang{6;7;6.5} \\
\sisetup{
angle-symbol-degree = d ,
angle-symbol-minute = m ,
angle-symbol-second = s
}
\ang{6;7;6.5}
}

\subsubsection{angle-symbol-over-decimal}
\ang{45.697} \\
\ang{6;7;6.5} \\
\ang[angle-symbol-over-decimal]{45.697} \\
\ang[angle-symbol-over-decimal]{6;7;6.5}\\
\ang[color=blue,angle-symbol-over-decimal]{6;7;6.5}
% \ang{10}\\
% \ang{12.3}\\
% \ang{4,5}\\
% \ang{1;2;3}\\
% \ang{;;1}\\
% \ang{+10;;}\\
% \ang{-0;1;}\\

\stepcounter{subsection}

\subsection{Using Units}
\subsubsection{inter-unit-product}
\unit{\farad\squared\lumen\candela} \\
\unit[inter-unit-product = \ensuremath{{}\cdot{}}]{\farad\squared\lumen\candela}\\
\unit[color=blue,inter-unit-product = \ensuremath{{}\cdot{}}]{\farad\squared\lumen\candela}

\subsubsection{per-mode, [display\textbar inline]-per-mode, per-symbol, fraction-command, bracket-unit-denominator}
\unit{\joule\per\mole\per\kelvin} \\
\unit{\metre\per\second\squared} \\
\unit[per-mode=fraction]{\joule\per\mole\per\kelvin} \\
\unit[per-mode=fraction]{\joule\raiseto{-1}\mole\per\kelvin} \\
\unit[per-mode=fraction]{\metre\per\second\squared}\\
\unit{\ampere\per\mole\second} \\
\unit[per-mode = power-positive-first]{\ampere\per\mole\second}

{
\sisetup{per-mode = symbol}%
\unit{\joule\per\mole\per\kelvin} \\
\unit{\metre\per\second\squared} \\
\unit[per-symbol = \ \text{div}\ ]{\joule\per\mole\per\kelvin} \\
\unit[bracket-unit-denominator = false]{\joule\per\mole\per\kelvin}\\
}
\unit[per-mode=repeated-symbol]{\joule\per\mole\per\kelvin}

{
\sisetup{per-mode = single-symbol}
\qty{10}{\per\metre} \\
\qty{20}{\metre\per\second} \\
\qty{30}{\joule\per\mole\per\kelvin}
}

{
\sisetup{
display-per-mode = fraction ,
inline-per-mode = symbol
}%
\( \unit{\joule\per\mole\per\kelvin} \)
\[ \unit{\joule\per\mole\per\kelvin} \]
\unit{\joule\per\mole\per\kelvin} \\
\(
\displaystyle
\unit{\joule\per\mole\per\kelvin}
\)
\[
\textstyle
\unit{\joule\per\mole\per\kelvin}
\]
\[
\textstyle
\unit[color=blue]{\joule\per\mole\per\kelvin}
\]
}

\newcommand{\myfraction}[2]{#1~\text{over }#2}
Other tests:\\
\unit[per-mode=single-symbol]{\meter\per\second\liter}\\
\unit[per-mode=fraction, fraction-command=\myfraction]{\joule\per\mole\per\kelvin}

\subsubsection{per-symbol-script-correction}
{
\sisetup{per-mode = symbol}%
\unit{\cm\cubed\per\gram} \\
\unit[per-symbol-script-correction = ]{\cm\cubed\per\gram} % todo spacing
}

\subsubsection{sticky-per}
\unit{\pascal\per\gray\henry} \\
\unit[sticky-per]{\pascal\per\gray\henry}

\subsubsection{qualifier-[mode\textbar phrase]}
\unit{\kilogram\polymer\squared\per\mole\catalyst\per\hour} \\
\unit[qualifier-mode = bracket]{\kilogram\polymer\squared\per\mole\catalyst\per\hour} \\
\unit[qualifier-mode = subscript]{\kilogram\polymer\squared\per\mole\catalyst\per\hour} \\
\unit[qualifier-mode = phrase, qualifier-phrase=\ ]{\kilogram\polymer\squared\per\mole\catalyst\per\hour} \\
\unit[qualifier-mode = combine]{\deci\bel\isotropic}\\
\unit{\kilogram\of{pol}\squared\per\mole\of{cat}\per\hour} \\
\unit[qualifier-mode = bracket]{\kilogram\of{pol}\squared\per\mole\of{cat}\per\hour} \\
\unit[qualifier-mode = combine]{\deci\bel\of{i}} \\
{
\sisetup{qualifier-mode = phrase, qualifier-phrase = \ }%
\unit{\kilogram\of{pol}\squared\per\mole\of{cat}\per\hour} \\
\sisetup{qualifier-phrase = \ \mbox{of}\ }%
\unit{\kilogram\of{pol}\squared\per\mole\of{cat}\per\hour}
}
%\unit[qualifier-mode = phrase]{\kilogram\polymer\squared\per\mole\catalyst\per\hour} \\
\unit[qualifier-mode = phrase, qualifier-phrase = { by }]{\kilogram\polymer\squared\per\mole\catalyst\per\hour}

% keep an eye on siunitx#854
\unit[qualifier-mode = bracket]{\kilogram\polymer\squared} \\  % v2 invalid; v3 kg(pol)^2
\unit[qualifier-mode = brackets]{\kilogram\polymer\squared} \\  % kg(pol)^2
\unit[qualifier-mode = phrase, qualifier-phrase = {\;}]{\kilogram\polymer\squared} \\ % v2 (kg pol)^2; v3 kg pol^2
\unit[qualifier-mode = space]{\kilogram\polymer\squared} \\ % v2 (kg pol)^2; v3 kg pol^2
\unit[qualifier-mode = combine]{\kilogram\polymer\squared} \\ % v2 invalid; v3 kgpol^2
\unit[qualifier-mode = text]{\kilogram\polymer\squared}  % kgpol^2

\subsubsection{power-half-as-sqrt}
\unit{\Hz\tothe{0.5}} \\
\unit[power-half-as-sqrt]{\Hz\tothe{0.5}}

\subsubsection{parse-units}
\qty{300}{\MHz} \\
\qty[parse-units = false]{300}{\MHz} \\
\unit{\meter\per\second} \\
\unit[parse-units=false]{\meter\per\second}
% Setting parse-units=false speeds things up by making many default assumptions.
% It's not clear which assumptions we miss.
% per-mode is the only one I know of.
% see siunitx#840

\subsubsection{forbid-literal-units}
\unit{m/s}\\
%\unit[forbid-literal-units]{m/s}\\  % raise an error
\unit[forbid-literal-units]{\meter\per\second}

\subsubsection{unit-font-command} % todo font
\unit{\lumen} \\
\unit[unit-font-command = \mathit]{\lumen}

\subsection{Quantities}
\subsubsection{allow-quantity-breaks}
\begin{minipage}{3cm}
% Gives an underfull hbox
X\hspace{2.4cm}\qty{10}{\metre} \\
X\hspace{2.1cm}\qty{10}{\metre} \\
\sisetup{allow-quantity-breaks}
X\hspace{2.4cm}\qty{10}{\metre} \\ % disagrees with texdoc? % todo spacing
X\hspace{2.1cm}\qty{10}{\metre} % todo spacing
\end{minipage}

\subsubsection{quantity-product} % todo spacing
\qty{2.67}{\farad} \\
\qty[quantity-product = \ ]{2.67}{\farad} \\
\qty[quantity-product = ]{2.67}{\farad}\\
\qty[quantity-product = \,]{2.67}{\farad}\\
\qty[quantity-product = \;]{2.67}{\farad}\\
\qty[quantity-product = $\times$]{2.67}{\farad}\\
\qty[color=blue,quantity-product = $\times$]{2.67}{\farad}\\
also\\
\qty[quantity-product = $\times$, separate-uncertainty]{3.7+-4.5}{\farad}


\subsubsection{prefix-mode, extract-mass-in-kilograms}
\qty{1e3}{\metre\second} \\
\qty[prefix-mode = combine-exponent]{1e3}{\metre\second} \\
\qty{10}{\kilo\gram\squared\deci\second} \\
\qty[prefix-mode = extract-exponent]{10}{\kilo\gram\squared\deci\second}\\
\qty[prefix-mode = extract-exponent]{7.5}{\gram} \\
{
\sisetup{extract-mass-in-kilograms = false}
\qty{10}{\kilo\gram\squared\deci\second} \\
\qty[prefix-mode = extract-exponent]{10}{\kilo\gram\squared\deci\second} \\
\qty[prefix-mode = extract-exponent]{7.5}{\gram}
}

\unit{\milli\litre\per\mole\deci\ampere} \\
%%\qty{10}{\kilo\gram\squared\deci\second} \\
\unit[prefix-mode = extract-exponent]{\milli\litre\per\mole\deci\ampere}\\
\unit[prefix-mode = extract-exponent]{\kilo\gram\squared\deci\second}\\
\unit{\mega\gram\squared\deci\second}\\
\unit[prefix-mode = extract-exponent]{\mega\gram\squared\deci\second}\\
\unit{\micro\gram\squared\deci\second}\\
\unit[prefix-mode = extract-exponent]{\micro\gram\squared\deci\second}\\
\unit{\per\mega\gram\squared\deci\second}\\
\unit[prefix-mode = extract-exponent]{\per\mega\gram\squared\deci\second}\\
\unit{\per\micro\gram\squared\deci\second}\\
\unit[prefix-mode = extract-exponent]{\per\micro\gram\squared\deci\second}
%%\qty[prefix-mode = extract-exponent]{10}{\kilo\gram\squared\deci\second}\\

\subsubsection{separate-uncertainty-units}
{
\sisetup{uncertainty-mode = separate}%
\qty{12.3(4)}{\kilo\gram} \\
\qty[separate-uncertainty-units = bracket]{12.3(4)}{\kilo\gram} \\
\qty[separate-uncertainty-units = repeat]{12.3(4)}{\kilo\gram}\\
\qty[separate-uncertainty-units = single]{12.3(4)}{\kilo\gram}
}

\subsection{Tabular Material}

\subsubsection{table-alignment-mode}

Parsing without aligning in an \texttt{S} column
\begin{center}
\begin{tabular}
{@{}S
S[table-alignment-mode = none]
@{}}
\toprule
{Decimal-centered} &
{Simple centering} \\
\midrule
12.345 & 12.345 \\
6,78 & 6,78 \\
-88.8(9) & -88.8(9) \\
4.5e3 & 4.5e3 \\
\bottomrule
\end{tabular}
\end{center}

\subsubsection{table-number-alignment}

Aligning the \texttt{S} column
\begin{center}
\sisetup{table-format = 2.4, table-alignment-mode = format}
\begin{tabular}{@{}
S[table-alignment-mode = marker]
S[table-number-alignment = center]
S[table-number-alignment = left]
S[table-number-alignment = right]
@{}}
\toprule
{Some Values} & {Some Values} & {Some Values} & {Some Values} \\
\midrule
2.3456 & 2.3456 & 2.3456 & 2.3456\\
34.2345 & 34.2345 & 34.2345 & 34.2345\\
56.7835 & 56.7835& 56.7835 & 56.7835\\
90.473 & 90.473 & 90.473 & 90.473\\
\bottomrule
\end{tabular}
\end{center}

\subsubsection{table-format}

Reserving space in \texttt{S} columns
\begin{center}
\sisetup{
table-alignment-mode = format,
table-number-alignment = center,
}
\begin{tabular}{@{}
S[table-format = 2.2]
S[table-format = 2.2, table-number-alignment = right]
S[table-format = 2.2(1)]
S[table-format = 2.2(1), uncertainty-mode = separate]
S[table-format = +2.2]
S[table-format = 2.2e1]
@{}}
\toprule
{Values}
& {Values}
& {Values}
& {Values}
& {Values}
& {Values} \\
\midrule
2.3 & 2.3 & 2.3(5) & 2.3(5) & 2.3 & 2.3e8 \\
34.23 & 34.23 & 34.23(4) & 34.23(4) & 34.23 & 34.23 \\
56.78 & 56.78 & 56.78(3) & 56.78(3) & -56.78 & 56.78e3 \\
3,76 & 3,76 & 3,76(2) & 3.76(2) & +-3.76 & e6 \\
\bottomrule
\end{tabular}
\end{center}

\subsubsection{table-model-setup}

\subsubsection{table-align-[comparator\textbar exponent\textbar uncertainty]}

The \texttt{table-align-exponent} option
\begin{center}
\sisetup{table-format = 1.3e2}
\begin{tabular}{@{}SS[table-align-exponent = false]@{}}
\toprule
{Header} & {Header} \\
\midrule
1.2e3 & 1.2e3 \\
1.234e56 & 1.234e56 \\
\bottomrule
\end{tabular}
\end{center}

The \texttt{table-align-uncertainty} option
\begin{center}
\sisetup{
uncertainty-mode = separate,
table-format = 1.3(1),
}
\begin{tabular}{@{}SS[table-align-uncertainty = false]@{}}
\toprule
{Header} & {Header} \\
\midrule
1.2(1) & 1.2(3) \\
1.234(5) & 1.234(5) \\
\bottomrule
\end{tabular}
\end{center}

The \texttt{table-align-comparator} option
\begin{center}
\sisetup{table-format = >2.2}
\begin{tabular}{@{}SS[table-align-comparator = false]@{}}
\toprule
{Header} & {Header} \\
\midrule
> 1.2 & > 1.2 \\
< 12.34 & < 12.34 \\
\bottomrule
\end{tabular}
\end{center}

Reserving space for comparators in \texttt{S} columns
\begin{center}
\sisetup{
table-number-alignment = center,
table-format=<<2.2e2,
}
\begin{tabular}{
S
S%[table-comparator = true]
S[table-align-comparator=false]}
\toprule
{Values} & {Values} & {Values} \\
\midrule
2 .3  & < 2.3e8    & < 2.3e8\\
34.23 & = 34.23    & = 34.23 \\
56.78 & >= 56.78e3 & >= 56.78e3\\
3,76  & \gg e6     & \gg e6 \\
\bottomrule
\end{tabular}
\end{center}

\subsubsection{table-align-text-[before\textbar after]} 

Closing notes up to text
\begin{center}
\newrobustcmd\NoteMark[1]{%
\textsuperscript{\emph{#1}}%
}
\sisetup{table-format = {\NoteMark{a}}2.4}
\begin{tabular}{@{}
S
S[table-align-text-before = false]
@{}}
\toprule
{Values} & {Values} \\
\midrule
2.3456 & 2.3456 \\
\NoteMark{a} 4.234 & \NoteMark{a} 4.234 \\
\NoteMark{b} .78 & \NoteMark{b} .78 \\
\NoteMark{d} 88 & \NoteMark{d} 88 \\
\bottomrule
\end{tabular}
\qquad
\sisetup{table-format = 2.4{\NoteMark{a}}}
%\sisetup{table-format = 2.4\NoteMark{a}}  % todo error
\begin{tabular}{@{}
S
S[table-align-text-after = false]
@{}}
\toprule
{Values} & {Values} \\
\midrule
2.3456 & 2.3456 \\
34.234 \NoteMark{a} & 34.234 \NoteMark{a} \\
56.78 \NoteMark{b} & 56.78 \NoteMark{b} \\
90.4 \NoteMark{c} & 90.4 \NoteMark{c} \\
88 \NoteMark{d} & 88 \NoteMark{d} \\
\bottomrule
\end{tabular}
\end{center}

\subsubsection{table-auto-round}

The \texttt{table-auto-round} option
\begin{center}
\sisetup{table-format = 1.3}
\begin{tabular}{@{}SS[table-auto-round]@{}}
\toprule
{Header} & {Header} \\
\midrule
1.2 & 1.2 \\
1.2345 & 1.2345 \\
\bottomrule
\end{tabular}
\end{center}

\subsubsection{parse-numbers}

Aligning without parsing
\begin{center}
\sisetup{
parse-numbers = false,
table-format = 3.3
}
\begin{tabular}{@{}
S
S[table-number-alignment = center]
S[table-number-alignment = right]
S[table-number-alignment = left]
@{}}
\toprule
{Some values}
& {Some values}
& {Some values}
& {Some values} \\
\midrule
2.35 & 2.35 & 2.35 & 2.35 \\
34.234 & 34.234 & 34.234 & 34.234 \\
56.783 & 56.783 & 56.783 & 56.783 \\
3,762 & 3,762 & 3,762 & 3.762 \\
\sqrt{2} & \sqrt{2} & \sqrt{2} & \sqrt{2} \\
\bottomrule
\end{tabular}
\end{center}

\subsubsection{drop-exponent}

The \texttt{drop-exponent} option
\begin{center}
\begin{tabular}{@{}
S[table-format = 1.1e1]
S[
drop-exponent = true,
exponent-mode = fixed,
fixed-exponent = 3,
table-format = 2.1,
]
@{}}
\toprule
{Header} & \multicolumn{1}{c@{}}{Header / \num[print-unity-mantissa = false]{e3}} \\
\midrule
1.2e3 & 1.2e3 \\
3e2 & 3e2 \\
1.0e4 & 1.0e4 \\
\bottomrule
\end{tabular}
\end{center}

\subsubsection{table-column-width, table-fixed-width}

Fixed-width columns
\begin{center}
\begin{tabular}
{@{}
S
S[table-column-width = 2cm]
@{}}
\toprule
{Flexible} &
{Fixed} \\
\midrule
1.23 & 1.23 \\
45.6 & 45.6 \\
\bottomrule
\end{tabular}
\end{center}

Right-aligning under a heading
\begin{center}
\newlength\mylength
\settowidth{\mylength}{Long header}
\sisetup{
table-alignment-mode = none ,
table-column-width = \mylength ,
table-number-alignment = right
}
\begin{tabular}{@{}S@{}}
\toprule
{Long header} \\
\midrule
12.33 \\
2 \\
1234 \\
\bottomrule
\end{tabular}
\end{center}

\subsubsection{table-text-alignment, table-alignment}

Aligning text in \texttt{S} columns
\begin{center}
\sisetup{table-format = 4.4}
\begin{tabular}{
S
S[table-text-alignment = left]
S[table-text-alignment = right]
}
\toprule
{Values}
& {Values}
& {Values} \\
\midrule
992.435 & 992.435 & 992.435 \\
7734.2344 & 7734.2344 & 7734.2344 \\
56.7834 & 56.7834 & 56.7834 \\
3,7462 & 3,7462 & 3,7462 \\
\bottomrule
\end{tabular}
\end{center}

\subsubsection{table-alignment}

\subsection{Locale Options}
\qty{1.234}{\metre}\\
\qty[locale = DE]{6.789}{\metre} % todo spacing

\subsection{Preamble-only options}
These are in a different file.  We check that trying to change the options has no effect here:

{
\sisetup{
	list-input-separator=:,
	product-input-separator=*,
	table-column-type=T
}
\numlist{1;2;3} \\
\numproduct{1x2x3} \\
Table check:\\
\begin{tabular}
{@{}S
S[table-alignment-mode = none]
@{}}
\toprule
{Decimal-centered} &
{Simple centering} \\
\midrule
12.345 & 12.345 \\
6,78 & 6,78 \\
-88.8(9) & -88.8(9) \\
4.5e3 & 4.5e3 \\
\bottomrule
\end{tabular}
}

% these are v2
%\subsubsection{input-product, input-quotient}
%\numproduct{1 x 2 x 3} \\
%\numproduct[input-product=*]{4 * 5 * 6} \\

%\section{Localisation}
%This is in a different file.

\setcounter{section}{9}

\section{Hints for using siunitx}

\setcounter{subsection}{1}
\subsection{Adjusting \textbackslash litre and \textbackslash liter}
{
\DeclareSIUnit\litre{SUCCESS}
\unit{\liter}
}

\setcounter{subsection}{4}

\subsection{Expanding content in tables}

Values as macros in \texttt{S} columns
\begin{center}
\newcommand*\myvaluea{1234}
\newcommand\myvalueb{1234}
\DeclareRobustCommand*\myvaluec{1234}
\protected\def\myvalued{1234}
\begin{tabular}{@{}S@{}}
\toprule
{Some Values} \\
\midrule
\myvaluea 8.8 \myvaluea \\ % Both expanded
\myvalueb 8.8 \myvalueb \\ % First expanded by TeX
% to numbers
%\myvaluec 8.8 \myvaluec \\ % First expanded by TeX  % todo error caused on next line
% but not to numbers!
\myvalued 8.8 \myvalued \\ % Neither expanded % todo output overlaps
{\myvaluea\ 8.8 \myvaluea} \\ % Neither expanded
\bottomrule
\end{tabular}
\end{center}

% \usepackage{xfp}
%Calculated values
%\begin{center}
%\newcommand{\valuea}{66.7012}
%\newcommand{\valueb}{66.0212}
%\newcommand{\valuec}{64.9026}
%\begin{tabular}{
%@{}
%S[table-format = 2.4]
%S[table-format = 3.4]
%@{}
%}
%\toprule
%{Value} & {Doubled} \\
%\midrule
%\valuea & \fpeval{\valuea * 2} \\
%\valueb & \fpeval{\valueb * 2} \\
%\valuec & \fpeval{\valuec * 2} \\
%\bottomrule
%\end{tabular}
%\end{center}

\setcounter{subsection}{10}

\subsection{Special considerations for the \texorpdfstring{\Backslash}{\textbackslash}\unit{\kWh} unit}

\unit{\kWh} \\
\unit{\kWh\per\metre}\\
{
\DeclareSIUnit\kWh{kWh}
\DeclareSIUnit\KWH{kWh}
\unit{\KWH\per\metre}\\
\unit{kWh.m^{-1}}\\
}
\unit{\candela\per\kWh}\\
\unit{\candela\per\kilo\watt\per\hour} \\
\unit[sticky-per]{\candela\per\kWh}

\subsection{Prefixes and small angles}

{
\DeclareSIUnit\arcsecond{as}
\qty{1e-3}{\arcsecond} \\
\qty[prefix-mode = combine-exponent]{1e-3}{\arcsecond}
}

%\setcounter{subsection}{15}
%
%\subsubsection{Symbolic `digits'}
%
%{
%\sisetup{input-digits = 0123456789\pi}%
%\num{4\pi e-7}\\ % fails
%\robustify\dots % fails
%\sisetup{input-digits = 0123456789\dots}%
%\qty{0,4066\dots}{\metre\squared}
%}

\setcounter{subsection}{16}

\subsection{Demonstrating prefixes}
\unit{\yotta\noop} \\
\qty[prefix-mode = extract-exponent, print-unity-mantissa = false]{1}{\yotta\noop} \\
a\unit{\noop}b

\end{document}
